\documentclass[
    aps,
    pre,
    twocolumn,
    superscriptaddress,
    nofootinbib
]{revtex4-2}

\usepackage[T2A]{fontenc}
\usepackage[utf8]{inputenc}
\usepackage[russian,english]{babel}
\usepackage{amsmath,amssymb}
\usepackage{graphicx}
\usepackage{booktabs}
\usepackage{siunitx}
\usepackage{xcolor}
\usepackage{hyperref}

\hypersetup{
    colorlinks=true,
    linkcolor=blue,
    citecolor=blue,
    urlcolor=blue
}

% -----------------------------------------------------------------------------
% Численные значения, которые удобно менять централизованно.
% TODO: перед финальной версией заменить на окончательные результаты fit.
% -----------------------------------------------------------------------------
\newcommand{\fitA}{5.520}
\newcommand{\fitB}{1.632}
\newcommand{\fitRMSE}{1.52\times10^{-3}}
\newcommand{\fitMAE}{9.76\times10^{-4}}
\newcommand{\fitN}{46}
\newcommand{\vdwTc}{1.002}
\newcommand{\vdwRhoc}{0.204}
\newcommand{\vdwPc}{0.0768}

\begin{document}

\title{Численная проверка области применимости уравнения Ван-дер-Ваальса\\
как эффективной аппроксимации изотерм Lennard--Jones-флюида}

\author{Стас Смекалин}
% TODO: указать официальное название организации и кафедры.
\affiliation{Московский физико-технический институт, Долгопрудный, Россия}
% TODO: при необходимости добавить научного руководителя отдельным \author.

\date{\today}

\begin{abstract}
Методом молекулярной динамики исследованы изотермы трехмерного флюида частиц,
взаимодействующих посредством потенциала Леннарда--Джонса. Расчеты выполнены в
периодической кубической ячейке в $NVT$-ансамбле в диапазоне приведенных
температур $0.7\leq T\leq1.5$ и на сетке плотностей, охватывающей разреженный газ,
область немонотонных изотерм и плотный флюид. По однофазной газовой части данных
восстановлены эффективные параметры уравнения Ван-дер-Ваальса,
$a=\fitA$ и $b=\fitB$; на использованных при подгонке точках получены
$\mathrm{RMSE}=\fitRMSE$ и $\mathrm{MAE}=\fitMAE$. Затем модель экстраполирована
за пределы области подгонки и сопоставлена с численными изотермами и приближенной
спинодалью, восстановленной по границам участков с
$(\partial P/\partial\rho)_T<0$. Уравнение Ван-дер-Ваальса правильно передает
температурный масштаб исчезновения механической неустойчивости, но заметно
искажает форму спинодали и плотностный масштаб переходной области. Зависимость
удельной потенциальной энергии от плотности при температурах выше критической
оказывается близкой к линейной, что согласуется с однородностью флюида и слабой
зависимостью его парной структуры от плотности в исследованном диапазоне.
Полученные результаты показывают, что двухпараметрическое уравнение
Ван-дер-Ваальса служит точной локальной аппроксимацией газовой ветви, но не дает
количественно корректного глобального описания Lennard--Jones-флюида.

% TODO: после окончательного анализа проверить все числа и при необходимости
% сократить аннотацию до 150--200 слов.
\end{abstract}

\maketitle

\section{Введение}

Связь между микроскопическим межчастичным взаимодействием и макроскопическим
уравнением состояния является одной из центральных задач статистической
физики. Потенциал Леннарда--Джонса представляет простейшую модель, в которой
одновременно присутствуют короткодействующее отталкивание и дальнее притяжение,
необходимые для возникновения газожидкостного перехода. Несмотря на
минимальность модели, ее термодинамика уже содержит критическую область,
метастабильные состояния и плотную жидкую фазу.

Уравнение Ван-дер-Ваальса является классическим среднеполевым приближением,
учитывающим исключенный объем частиц и среднее притяжение двумя постоянными
$a$ и $b$. Оно качественно воспроизводит существование критической точки и
немонотонной докритической изотермы, однако не следует непосредственно из
потенциала Леннарда--Джонса и не обязано количественно описывать его во всем
диапазоне состояний. Поэтому параметры $a$ и $b$ в настоящей работе
рассматриваются как эффективные параметры, определенные для явно заданной
области температур и плотностей.

Цель работы состоит в численном определении изотерм Lennard--Jones-флюида,
локальной подгонке уравнения Ван-дер-Ваальса по однофазной газовой области и
проверке того, какие свойства численной модели сохраняются при экстраполяции
этой аппроксимации к фазово-неустойчивым и плотным состояниям. Особое внимание
уделяется сравнению спинодалей и анализу удельной потенциальной энергии как
дополнительной характеристики структурной перестройки флюида.

% TODO: добавить 3--5 ссылок на стандартные работы по LJ-флюиду, уравнениям
% состояния и прямому определению спинодали по молекулярной динамике.

\section{Модель и численный метод}

\subsection{Потенциал взаимодействия и приведенные единицы}

Рассматривалась система из $N$ одинаковых частиц массы $m$, взаимодействующих
посредством парного потенциала Леннарда--Джонса
\begin{equation}
    u(r)=4\varepsilon
    \left[
        \left(\frac{\sigma}{r}\right)^{12}
        -
        \left(\frac{\sigma}{r}\right)^6
    \right].
    \label{eq:lj}
\end{equation}
Вычисления проводились в приведенных единицах
\begin{equation}
    \sigma=\varepsilon=m=k_{\mathrm B}=1,
\end{equation}
так что температура, давление, плотность и энергия далее приводятся в
соответствующих безразмерных LJ-единицах. Потенциал обрезался на расстоянии
$r_c=2.5\sigma$.

% TODO: точно указать способ обрезания: простой cutoff, shifted potential,
% shifted force и наличие или отсутствие tail correction.

Частицы находились в кубической ячейке объема $V=L^3$ с периодическими
граничными условиями. Средняя численная плотность задавалась как
\begin{equation}
    \rho=\frac{N}{V}.
\end{equation}

\subsection{Молекулярная динамика}

Динамика интегрировалась в $NVT$-ансамбле с использованием термостата Ланжевена.
Для каждой пары $(T,\rho)$ выполнялось несколько независимых прогонов с разными
случайными начальными условиями. Каждый прогон состоял из участка релаксации,
данные которого отбрасывались, и последующего участка измерения, на котором
усреднялись температура, давление, потенциальная и кинетическая энергии.

% TODO: вставить окончательные параметры расчета:
% N = ...;
% dt = ...;
% коэффициент трения = ...;
% equil_steps = ...;
% prod_steps = ...;
% sample_interval = ...;
% число seed = ...;
% полный список температур и плотностей.

Начальные координаты формировались без перекрытий частиц, после чего начальные
скорости выбирались из распределения Максвелла при заданной температуре.
Расчеты выполнялись с помощью пакета OpenMM на графическом ускорителе.

% TODO: указать версию OpenMM, тип GPU и точность CUDA, если это требуется в
% отчете. Подробности облачной инфраструктуры не включать.

\subsection{Давление и статистические погрешности}

Давление в периодической системе определялось по вириальной формуле
\begin{equation}
    P=\rho T+\frac{1}{3V}
    \left\langle
        \sum_{i<j}\mathbf r_{ij}\cdot\mathbf F_{ij}
    \right\rangle,
    \label{eq:virial-pressure}
\end{equation}
где $\mathbf r_{ij}$ --- вектор минимального образа между частицами,
а $\mathbf F_{ij}$ --- действующая между ними сила. Первый член в
Eq.~\eqref{eq:virial-pressure} является идеальногазовым вкладом, тогда как
вириальный член учитывает притяжение и отталкивание. Для разреженного газа
вириальная поправка стремится к нулю и выполняется $P\to\rho T$.

Для каждой пары $(T,\rho)$ результаты независимых прогонов усреднялись по seed.
Погрешность точки оценивалась по разбросу независимых средних значений давления.
На графиках показаны средние значения и соответствующие интервалы
неопределенности.

% TODO: точно описать используемую оценку погрешности: стандартное отклонение,
% стандартная ошибка среднего, доверительный интервал или комбинация временной
% и межсидовой дисперсии. Привести формулу, совпадающую с кодом.

\section{Численные изотермы}

На рис.~\ref{fig:isotherms} представлены полученные зависимости давления от
плотности для температур от $T=0.7$ до $T=1.5$. При малых плотностях все кривые
приближаются к идеальногазовому пределу $P=\rho T$. При увеличении плотности
притягивающая часть взаимодействия уменьшает давление относительно
идеальногазового значения. На докритических изотермах возникает немонотонный
участок, на котором давление уменьшается при росте плотности и в части диапазона
может принимать отрицательные значения. При дальнейшем сжатии доминирует
короткодействующее отталкивание, и давление быстро возрастает на плотной ветви.

\begin{figure*}[t]
    \centering
    % TODO: заменить placeholder на реальный файл, например:
    % \includegraphics[width=0.92\textwidth]{figures/eos_isotherms.pdf}
    \fbox{\parbox[c][0.28\textheight][c]{0.9\textwidth}{\centering
    Вставить семейство изотерм $P(\rho)$ для всех температур с погрешностями.}}
    \caption{Изотермы Lennard--Jones-флюида, полученные методом молекулярной
    динамики. Показаны средние значения давления по независимым прогонам;
    вертикальные отрезки соответствуют статистической погрешности. Пунктиром
    может быть показан идеальногазовый предел $P=\rho T$ для нескольких
    характерных температур.}
    \label{fig:isotherms}
\end{figure*}

Отрицательный наклон конечносистемной изотермы не следует интерпретировать как
устойчивую макроскопическую ветвь равновесного вещества. Для однородной фазы
механическая устойчивость требует
\begin{equation}
    \left(\frac{\partial P}{\partial\rho}\right)_T>0.
\end{equation}
Нарушение этого условия соответствует области механической неустойчивости
однородного состояния. В конечной периодической $NVT$-системе наблюдаемая форма
кривой дополнительно зависит от времени релаксации, размера ячейки,
метастабильности и морфологии неоднородного состояния. Поэтому далее
немонотонные участки используются для построения приближенной численной
спинодали, но не отождествляются с точной термодинамической кривой сосуществования.

% TODO: после финализации графика добавить 2--3 численных наблюдения:
% при каких T немонотонность видна;
% где впервые появляется P<0;
% при какой T немонотонность исчезает в пределах погрешности.

\section{Удельная потенциальная энергия}

Зависимость удельной потенциальной энергии $U/N$ от плотности показана на
рис.~\ref{fig:energy}. При $\rho\to0$ среднее число взаимодействующих соседей
стремится к нулю и $U/N\to0$. С увеличением плотности потенциальная энергия
становится отрицательной вследствие притяжения между частицами. На низких
температурах форма кривой заметно меняется в области структурной перестройки и
образования плотных связанных областей.

\begin{figure*}[t]
    \centering
    % TODO: заменить на реальный график.
    % \includegraphics[width=0.92\textwidth]{figures/potential_energy.pdf}
    \fbox{\parbox[c][0.28\textheight][c]{0.9\textwidth}{\centering
    Вставить семейство зависимостей $U/N$ от $\rho$.}}
    \caption{Средняя потенциальная энергия на одну частицу как функция
    плотности. При температурах выше критической зависимость близка к линейной,
    тогда как на докритических изотермах наблюдается выраженная структурная
    перестройка.}
    \label{fig:energy}
\end{figure*}

Для однородного изотропного флюида средняя потенциальная энергия выражается
через радиальную функцию распределения $g(r;\rho,T)$:
\begin{equation}
    \frac{U}{N}
    =2\pi\rho\int_0^{r_c}
    u(r)g(r;\rho,T)r^2\,dr.
    \label{eq:energy-gr}
\end{equation}
Если при фиксированной температуре форма $g(r;\rho,T)$ слабо изменяется с
плотностью, интеграл в Eq.~\eqref{eq:energy-gr} остается приблизительно
постоянным и возникает линейная зависимость
\begin{equation}
    \frac{U}{N}\simeq C(T)\rho.
    \label{eq:linear-energy}
\end{equation}
Именно такое поведение наблюдается для температур выше критической в
исследованном диапазоне плотностей. Оно согласуется с отсутствием
макроскопического разделения фаз и со сравнительно слабым изменением локальной
парной структуры при сжатии. Нарушение линейности на более низких температурах
служит независимым качественным признаком структурной перестройки и согласуется
с критической температурой, оцененной по исчезновению немонотонности изотерм.

% TODO: выполнить и привести линейные fit для нескольких T>T_c:
% коэффициенты наклона, R^2 или относительные остатки.
% Не утверждать, что линейность сама по себе строго определяет T_c; использовать
% ее как дополнительное согласующееся свидетельство.

\section{Эффективное уравнение Ван-дер-Ваальса}

\subsection{Область подгонки}

Численные данные сопоставлялись с уравнением
\begin{equation}
    P_{\mathrm{VdW}}(\rho,T)
    =\frac{\rho T}{1-b\rho}-a\rho^2.
    \label{eq:vdw}
\end{equation}
Параметры $a$ и $b$ определялись только по точкам однофазной газовой и
умеренно-плотной области. Такое ограничение необходимо, поскольку внутри
двухфазной области и на плотной жидкой ветви среднеполевое уравнение
Eq.~\eqref{eq:vdw} не обязано воспроизводить структуру конечносистемных
MD-изотерм. В текущей обработке использованы точки с
$T\geq1.1$ и $\rho\leq0.20$.

% TODO: проверить и окончательно зафиксировать область fit. Если границы были
% другими, заменить предложение выше и указать число точек.
% TODO: описать точную целевую функцию: обычные наименьшие квадраты или fit с
% весами 1/sigma_P^2.

Полученные значения приведены в табл.~\ref{tab:fit}. Малые значения RMSE и MAE
показывают, что в выбранной области двухпараметрическая модель локально хорошо
описывает численные давления.

\begin{table}[b]
    \caption{Предварительные результаты подгонки уравнения
    Ван-дер-Ваальса.}
    \label{tab:fit}
    \begin{ruledtabular}
        \begin{tabular}{lc}
            Величина & Значение \\
            \hline
            $a$ & $\fitA$ \\
            $b$ & $\fitB$ \\
            $\mathrm{RMSE}(P)$ & $\fitRMSE$ \\
            $\mathrm{MAE}(P)$ & $\fitMAE$ \\
            Число точек & $\fitN$ \\
        \end{tabular}
    \end{ruledtabular}
\end{table}

\subsection{Критические параметры модели}

Для уравнения Ван-дер-Ваальса критические параметры выражаются через найденные
коэффициенты:
\begin{equation}
    \rho_c^{\mathrm{VdW}}=\frac{1}{3b},
    \qquad
    T_c^{\mathrm{VdW}}=\frac{8a}{27b},
    \qquad
    P_c^{\mathrm{VdW}}=\frac{a}{27b^2}.
    \label{eq:vdw-critical}
\end{equation}
Для текущей подгонки это дает
\begin{equation}
    \rho_c^{\mathrm{VdW}}\approx\vdwRhoc,
    \qquad
    T_c^{\mathrm{VdW}}\approx\vdwTc,
    \qquad
    P_c^{\mathrm{VdW}}\approx\vdwPc.
\end{equation}

Температурный масштаб критической области оказывается близок к численной оценке,
полученной по исчезновению немонотонности MD-изотерм. При этом критическая
плотность и форма переходной области воспроизводятся значительно хуже. Такое
сочетание является важным результатом: хороший локальный fit газовой ветви
может давать разумную оценку характерной температуры, не обеспечивая
количественно верного описания всей фазовой диаграммы.

% TODO: вставить численную оценку T_c по MD и ее неопределенность/разрешение по
% температурной сетке. Явно сравнить ее с T_c^VdW.

\subsection{Сравнение изотерм и остатки}

На рис.~\ref{fig:vdw-comparison} показано сравнение численных изотерм с
Eq.~\eqref{eq:vdw}. В области подгонки различия малы и не имеют выраженного
систематического характера. При переходе к меньшим температурам и большим
плотностям отклонения становятся структурированными: модель неверно описывает
положение экстремумов, ширину немонотонной области и рост давления на плотной
ветви.

\begin{figure*}[t]
    \centering
    % TODO: заменить на реальную сравнительную фигуру.
    % \includegraphics[width=0.92\textwidth]{figures/vdw_comparison.pdf}
    \fbox{\parbox[c][0.28\textheight][c]{0.9\textwidth}{\centering
    Вставить сравнение MD-точек и изотерм Ван-дер-Ваальса для нескольких
    характерных температур.}}
    \caption{Сравнение изотерм, полученных методом молекулярной динамики, с
    эффективным уравнением Ван-дер-Ваальса. Следует показать как минимум одну
    температуру из области fit, одну околокритическую и одну докритическую.}
    \label{fig:vdw-comparison}
\end{figure*}

Для количественного представления расхождения использовался остаток
\begin{equation}
    \Delta P(\rho,T)
    =P_{\mathrm{MD}}(\rho,T)-P_{\mathrm{VdW}}(\rho,T).
    \label{eq:residual}
\end{equation}
Его зависимость от плотности позволяет отличить статистический разброс от
систематического нарушения формы уравнения состояния.

% TODO: вставить график остатков или решить, достаточно ли сравнительного
% графика. Если остатки не добавляются в отчет, удалить этот абзац.
% TODO: привести RMSE отдельно внутри и вне области fit, если это уже посчитано.

\section{Численная оценка спинодали}

Спинодаль однородной фазы определяется условием потери механической устойчивости
\begin{equation}
    \left(\frac{\partial P}{\partial\rho}\right)_T=0.
    \label{eq:spinodal-condition}
\end{equation}
Между двумя ветвями спинодали выполняется
$(\partial P/\partial\rho)_T<0$. В настоящей работе использовалась дискретная
численная оценка: для каждой температуры выбирались измеренные точки,
ограничивающие связный немонотонный участок изотермы. Такой способ не вводит
точности, превосходящей разрешение исходной сетки плотностей, и позволяет
непосредственно связать спинодаль с наблюдаемыми MD-данными.

Если немонотонный участок отсутствовал либо его границы не определялись надежно
из-за статистических погрешностей, соответствующая температура не включалась в
автоматическое построение. Полученная кривая далее называется приближенной
спинодалью конечносистемных MD-изотерм.

Для уравнения Ван-дер-Ваальса условие
Eq.~\eqref{eq:spinodal-condition} принимает вид
\begin{equation}
    \frac{T}{(1-b\rho)^2}-2a\rho=0,
    \label{eq:vdw-spinodal}
\end{equation}
или, эквивалентно,
\begin{equation}
    T_{\mathrm{sp}}^{\mathrm{VdW}}(\rho)
    =2a\rho(1-b\rho)^2.
    \label{eq:vdw-spinodal-explicit}
\end{equation}

\begin{figure}[t]
    \centering
    % TODO: заменить на реальную диаграмму.
    % \includegraphics[width=0.95\columnwidth]{figures/spinodal.pdf}
    \fbox{\parbox[c][0.25\textheight][c]{0.9\columnwidth}{\centering
    Вставить сравнение численной спинодали и кривой Ван-дер-Ваальса в
    координатах $(\rho,T)$.}}
    \caption{Приближенная спинодаль, восстановленная по границам
    немонотонных участков MD-изотерм, и спинодаль эффективного уравнения
    Ван-дер-Ваальса.}
    \label{fig:spinodal}
\end{figure}

Сравнение на рис.~\ref{fig:spinodal} показывает, что уравнение
Ван-дер-Ваальса правильно воспроизводит общий температурный масштаб исчезновения
немонотонной области. В то же время форма двух ветвей и соответствующие им
плотности заметно отличаются от численных данных. Следовательно, совпадение по
критической температуре не означает правильного воспроизведения геометрии
области метастабильности.

% TODO: добавить конкретные числа: диапазон T, в котором найдена MD-спинодаль;
% положение ее верхней точки; максимальное/характерное расхождение по rho.
% TODO: если будет построена бинодаль, добавить отдельный подраздел и четко
% отделить ее от спинодали. Не смешивать эти два понятия.

\section{Обсуждение}

Полученные результаты разделяют локальную и глобальную применимость уравнения
Ван-дер-Ваальса. В газовой однофазной области параметры $a$ и $b$ обеспечивают
малую ошибку давления. Здесь среднеполевое описание успешно учитывает два
главных эффекта взаимодействия: уменьшение доступного объема и снижение
давления вследствие притяжения.

При переходе к докритическим и плотным состояниям предпосылки модели становятся
недостаточными. Реальный Lennard--Jones-флюид обладает зависящей от плотности и
температуры парной структурой, тогда как Eq.~\eqref{eq:vdw} кодирует ее двумя
постоянными коэффициентами. Поэтому модель не может одновременно точно
воспроизвести газовую ветвь, форму спинодали и быстрое возрастание давления в
плотной фазе.

Близость температурного масштаба критической области при различии плотностного
масштаба показывает, что эффективная аппроксимация сохраняет некоторую
интегральную информацию о балансе притяжения и теплового движения. Однако она
не восстанавливает детальную структуру корреляций, определяющую геометрию
фазовой диаграммы.

Интерпретация численной спинодали ограничена конечным размером системы и
конечным временем наблюдения. Немонотонная ветвь может зависеть от
метастабильности, скорости структурной релаксации и конкретной морфологии
неоднородного состояния в периодической ячейке. Поэтому построенная кривая
служит операционным показателем границ наблюдаемой механической неустойчивости,
а не точным термодинамическим пределом бесконечной системы.

% TODO: при наличии данных добавить проверку зависимости от N, числа шагов или
% seed. Если таких проверок нет, не обещать их и оставить ограничение явно.
% TODO: после готовности ролика можно добавить одно предложение о визуализации
% неравновесного изотермического сжатия как иллюстрации микроскопической
% перестройки, но не использовать ролик как количественное доказательство.

\section{Заключение}

Методом молекулярной динамики получено семейство изотерм трехмерного
Lennard--Jones-флюида в широком диапазоне температур и плотностей. При малых
плотностях результаты переходят в идеальногазовый предел, тогда как на
докритических изотермах наблюдаются немонотонные участки и состояния с
отрицательным давлением. При дальнейшем сжатии давление быстро возрастает из-за
короткодействующего отталкивания.

По газовой однофазной части данных восстановлены эффективные параметры
уравнения Ван-дер-Ваальса. В области подгонки модель описывает давление с малой
ошибкой, однако ее экстраполяция обнаруживает систематические расхождения с
MD-данными. Спинодаль Ван-дер-Ваальса близка к численной по температурному
масштабу, но существенно отличается по форме и положению ветвей в пространстве
плотностей.

Почти линейная зависимость $U/N$ от $\rho$ при температурах выше критической
согласуется с однородностью флюида и слабым изменением его парной структуры при
сжатии. Нарушение этой линейности при более низких температурах отражает
структурную перестройку и дополняет выводы, полученные из изотерм давления.

Таким образом, уравнение Ван-дер-Ваальса является содержательной и точной
локальной аппроксимацией газовой области Lennard--Jones-флюида, но не дает
универсального количественного описания фазово-неустойчивых и плотных состояний.
Именно различие между хорошим локальным fit и неудачной глобальной
экстраполяцией определяет границу применимости этой простой эффективной модели.

% TODO: перед финальной версией привести здесь 3--4 окончательных численных
% результата: a, b, T_c^MD, T_c^VdW и один показатель качества fit.

\begin{acknowledgments}
% TODO: указать научного руководителя, преподавателя, вычислительные ресурсы и
% другие благодарности. Если раздел не нужен, удалить целиком.
Автор благодарит [имя] за обсуждение постановки задачи и замечания по работе.
\end{acknowledgments}

\appendix

\section{Дополнительные параметры расчета}

% TODO: этот раздел можно удалить, если все параметры помещаются в основном
% тексте, либо оформить здесь компактную таблицу полной сетки эксперимента.
\begin{table}[h]
    \caption{Параметры численного эксперимента.}
    \label{tab:simulation-parameters}
    \begin{ruledtabular}
        \begin{tabular}{lc}
            Параметр & Значение \\
            \hline
            Число частиц $N$ & [вставить] \\
            Шаг интегрирования $\Delta t$ & [вставить] \\
            Радиус обрезания $r_c$ & $2.5\sigma$ \\
            Температуры & $0.7\ldots1.5$ \\
            Плотности & [вставить диапазон и сетку] \\
            Число независимых seed & [вставить] \\
            Шаги релаксации & [вставить] \\
            Шаги измерения & [вставить] \\
        \end{tabular}
    \end{ruledtabular}
\end{table}

\begin{thebibliography}{99}

\bibitem{AllenTildesley}
M. P. Allen and D. J. Tildesley,
\textit{Computer Simulation of Liquids}, 2nd ed.
(Oxford University Press, Oxford, 2017).

\bibitem{FrenkelSmit}
D. Frenkel and B. Smit,
\textit{Understanding Molecular Simulation}, 2nd ed.
(Academic Press, San Diego, 2002).

\bibitem{HansenMcDonald}
J.-P. Hansen and I. R. McDonald,
\textit{Theory of Simple Liquids}, 4th ed.
(Academic Press, Oxford, 2013).

\bibitem{OpenMM}
P. Eastman \textit{et al.},
OpenMM 7: Rapid development of high performance algorithms for molecular dynamics,
PLoS Comput. Biol. \textbf{13}, e1005659 (2017).

% TODO: добавить источник по фазовой диаграмме LJ/LJTS с тем же cutoff и тем же
% способом обрезания потенциала.
% TODO: добавить источник по уравнению состояния/критической точке LJ-флюида.
% TODO: добавить источник по спинодали или механической устойчивости.

\end{thebibliography}

\end{document}